%%
%% This is file `coppe-quickref.tex',
%% generated with the docstrip utility.
%%
%% The original source files were:
%%
%% ufrj-coppe.dtx  (with options: `quickrefcoppe')
%% 
%% This is a sample monograph which illustrates the use of the `ufrj' document
%% class with the `ufrj-coppe' unit style and its biblatex/biber bibliography
%% style. It is generated from ufrj-coppe.dtx by docstrip.
%% 
%% Copyright (C) 2008-2026 Vicente Helano F. Batista, George O. Ainsworth Jr.,
%% Geraldo Xexéo and Eduardo Mangeli. Released under the GNU General Public
%% License version 3 (see the COPYING.txt file).
%% 
%% coppe-quickref.tex -- CoppeTeX quick reference, COPPE unit style, in English.
%%
%% GERADO de ufrj-coppe.dtx. Nao edite este arquivo: edite o modulo
%% `quickrefcoppe' do ufrj-coppe.dtx e rode `pdflatex ufrj-coppe.ins'.
\documentclass[a4paper,10pt]{article}
\usepackage[T1]{fontenc}
\usepackage[utf8]{inputenc}
\usepackage{lmodern}
\usepackage[margin=2cm]{geometry}
\usepackage{longtable}
\usepackage{booktabs}
\usepackage{array}
\usepackage{xcolor}
\usepackage[colorlinks=true,allbordercolors=white]{hyperref}

\newcommand\ufrjversion{v5.0}

\newcolumntype{C}{>{\ttfamily\raggedright\arraybackslash}p{0.17\textwidth}}
\newcolumntype{D}{>{\raggedright\arraybackslash}p{0.44\textwidth}}
\newcolumntype{T}{>{\raggedright\arraybackslash}p{0.30\textwidth}}

\setlength\parindent{0pt}

\begin{document}

\begin{center}
  {\Large\bfseries CoppeTeX quick reference --- COPPE}\\[2pt]
  {\normalsize Unit style \texttt{ufrj-coppe}, \ufrjversion}\\[6pt]
  \begin{minipage}{0.92\textwidth}\small
  This sheet is the COPPE half. The other half --- every option, command and
  environment of the class, which is the Manual UFRJ/SiBI and names no unit ---
  is \texttt{ufrj-quickref.pdf}; the full manual of the style is
  \texttt{ufrj-coppe.pdf}, in Portuguese.

  \medskip
  A COPPE thesis is the class plus this style:

  \medskip
  \verb|\documentclass[dsc]{ufrj}|\\
  \verb|\usepackage{ufrj-coppe}|

  \medskip
  The style brings the institute's name on the cover and in the reference above
  each abstract, its thirteen Programas, the COPPE logo to the right of UFRJ's,
  the sentence that says to whom the work was presented, the name of the degree
  (``Mestre/Doutor em Ciências''), the city (Rio de Janeiro, RJ) and the norm
  quoted in the colophon. Everything else is the class.
  \end{minipage}
\end{center}

\vspace{6pt}

\begin{longtable}{@{}CTD@{}}
\toprule
\normalfont\bfseries \textbackslash department & \normalfont\bfseries Programa
& \normalfont\bfseries In English \\
\midrule
\endfirsthead
\toprule
\normalfont\bfseries \textbackslash department & \normalfont\bfseries Programa
& \normalfont\bfseries In English \\
\midrule
\endhead
\bottomrule
\endfoot
PEB  & Engenharia Biomédica & Biomedical Engineering \\
PEC  & Engenharia Civil & Civil Engineering \\
PEE  & Engenharia Elétrica & Electrical Engineering \\
PEM  & Engenharia Mecânica & Mechanical Engineering \\
PEMM & Engenharia Metalúrgica e de Materiais & Metallurgical and Materials Engineering \\
PEN  & Engenharia Nuclear & Nuclear Engineering \\
PENO & Engenharia Oceânica & Ocean Engineering \\
PENT & Engenharia de Nanotecnologia & Nanotechnology Engineering \\
PEP  & Engenharia de Produção & Production Engineering \\
PEQ  & Engenharia Química & Chemical Engineering \\
PESC & Engenharia de Sistemas e Computação & Systems Engineering and Computer Science \\
PET  & Engenharia de Transportes & Transportation Engineering \\
PPE  & Planejamento Energético & Energy Planning \\
\end{longtable}

\vspace{4pt}
\begin{center}
\begin{minipage}{0.92\textwidth}\small
\textbf{An acronym the style does not declare is a class error}, and not a
cover without the Programa. A Programa created after this release can be
declared in the preamble, after the style:

\medskip
\verb|\ufrjdeclareprogram{PEX}{Engenharia de Exemplo}{Example Engineering}|

\medskip
\textbf{The COPPE norm.} The colophon says the work follows the
\emph{Norma para a Elaboração Gráfica de Teses e Dissertações da COPPE/UFRJ}.
The norm decides only where the Manual does not oblige --- it adds the lists of
\emph{quadros}, of programs and of algorithms, and keeps the three elements at
the head of the abstract sheet, which \verb|\setupabstracts| turns on. Where
the two disagree, the Manual wins.

\medskip
\textbf{A thesis begun before v4.1} compiles as it is: \texttt{coppe.cls} is a
compatibility class that loads \texttt{ufrj} with the same options, loads this
style and warns once. The old public names --- \verb|\coppestring| and its
brothers, \verb|\usecoppelanguage|, \verb|\newcoppefloat|,
\verb|\coppetexfinalpage| and the six colophon pieces \verb|\coppefinal...| ---
go on working, and a \verb|\renewcommand| by the old name still changes the
colophon. It relaxes no rule of the class.

\medskip
Worked examples: \texttt{coppe-min-exemplo.tex}, only what is mandatory, and
\texttt{coppe-max-exemplo.tex}, everything the class composes.\quad
Problems: \url{https://github.com/COPPE-UFRJ/CoppeTeX/issues}
\end{minipage}
\end{center}

\end{document}
%% 
%%
%% End of file `coppe-quickref.tex'.
